\section*{Annotated plan for a proof companion}

\textbf{Purpose.}
Translate the arguments in the Lean development into mathematical prose that can be
read without following tactic scripts. The exposition should explain how each conclusion
follows from the definitions and intermediate results actually used in its formal proof.
The translations below cover the Cook--Levin proof route and sparse-counter signed
powers. The other entries are proposed scope, not claims of completed translations.

\subsection*{A common form for each proof}
State the exact theorem, its model and hypotheses. Introduce the objects needed to read
its argument. Then explain the intermediate claims, their proofs and how they compose,
with links to the corresponding Lean declarations. Make uses of choice, finiteness,
encoding conventions and quantitative bounds explicit where they matter. Cite the
mathematical sources and credit earlier formalizations. A named imported lemma can be
used as a lemma, with its contract stated and a link for following the argument further.
\choice{How far down should each translation go? The sample explains the tableau
induction and its one-hot argument, while treating the machine's local step-equation
lemma and previously certified list operations as named prerequisites.}

\subsection*{1. Computation and reductions}
Explain the machine and encoding contracts needed by later proofs; translate the
composition and iteration arguments behind resource certificates. Use Cook--Levin as
the first complete route from a mathematical construction to a class-level result.
Subsequent sections could translate the universal simulation, selected complete-problem
reductions and the time hierarchy, stating the precise overhead in each case.
\choice{We can begin with Cook--Levin and introduce foundational lemmas on demand,
or place a short model chapter first.}

\subsection*{2. Space, alternation and interaction}
Translate the recursive reachability argument for the polynomial-space Savitch result,
the two containments giving $\cc{AP}=\cc{PSPACE}$, and the game evaluation behind
$\cc{IP}\subseteq\cc{PSPACE}$. Explain the actual stored state and recursion bounds.
Follow any characterization theorem used by these proofs back to its machine interface.
\choice{Keep these as separate proofs with cross-references to shared space lemmas,
or group them around the common problem of evaluating large computation trees.}

\clearpage
\subsection*{3. Randomness, counting and cryptography}
Translate the covering argument in Sipser--G\'acs--Lautemann; the isolation, parity and
counting steps in Toda; and the hybrid and reduction arguments used for PRGs and
adaptive PRFs. State the oracle, adversary and probability spaces used by the formal
proofs. Follow each claimed efficiency guarantee to its certificate.
\choice{Which proof should follow Cook--Levin? Toda would exercise quantitative
probabilistic reasoning; a PRG/PRF proof would exercise the security interfaces.}

\subsection*{4. Circuits, logic and combinatorial bounds}
The sparse-counter signed-power argument is translated below. Further chapters could
include Fagin's theorem, determinant circuits in $\cc{VP}$,
resolution size--width, and expander-walk hitting. For each, explain the construction or
inequality chain in the actual Lean development. Keep formal polynomials distinct from
polynomial functions, and proof certificates distinct from runtime analyses.
\choice{Choose a small initial set for depth. The remaining results can stay in this
coverage plan until their arguments have been translated and reviewed.}

\textbf{Priority gaps.}
First-order reductions and their restricted projection forms are a priority, building
on the existing interpretations and formula pullbacks. Circuit Value is currently proved
in P; its logspace P-completeness theorem remains to be formalized. The NC/AC/TC
hierarchies are nonuniform; uniform-family and machine-characterization theorems are
separate work. Likewise, the verifier definitions of \#P and GapP still need their
nondeterministic-machine characterizations. These are development targets, not translated
proofs already present in the library.

\subsection*{5. The PCP theorem}
Translate the route through gap amplification and composition to
$\cc{NP}=\cc{PCP}(O(\log n),O(1))$, including the parameter calculations and the
certificates for construction and verification. Preserve the proof's dependencies
through explicit lemma references. Explain where the large fixed constants arise.
\choice{This is likely a sequence of proof sections with its own dependency map,
rather than one overview chapter.}

\subsection*{Review and maintenance}
An independent AI reviewer compares each translation with the cited Lean definitions,
statements and proof bodies at a recorded revision. The review checks hypotheses,
quantifier order, intermediate claims, proof steps and resource bounds, and records any
omission as an explicit prerequisite or an unresolved gap. AI review is evidence of
inspection; Lean does not check equivalence between the prose and the formal proof.
The author continues to control wording, organization and selection.

Every AI-authored PR must synchronize the affected exposition with its final Lean
changes and obtain a fresh independent review of the changed translation. A proof
refactor may change the argument even when the theorem statement stays fixed.
New mathematical developments need corresponding exposition; renamed or moved
results need repaired links. A PR with no mathematical or exposition impact must say
so explicitly. Its review must not claim coverage that has not been written.

\textbf{Coverage and navigation.}
A small index will associate translated arguments with their source modules and
revision, named prerequisites, and review status. It should make the boundary between
reviewed exposition and proposed chapters visible. This preserves the old blueprint's
links from readable mathematics to Lean, while giving the proofs themselves room to be read.
